\section{Reward Shaping 与价值稳定}

\subsection{价值函数稳定的微调技巧}

在 off-policy 更新中，Q 函数预测 drift 会严重影响 actor。实践中常见的 stabilizer 包括：

\begin{itemize}
    \item \textbf{Target smoothing}：在 TD target 中使用 twin Q average + Polyak average smoothing coefficient $\tau$。
    \item \textbf{Reward clipping}：将 reward 限制在 $[-1,1]$，避免爆炸性 gradient。
    \item \textbf{Value normalization}：在 batch_level 对 Q 值做人口归一化，减少 scale drift。
\end{itemize}

\subsection{Shaping 的设计原则}

Reward shaping 需要遵从 potential-based shaping theorem，确保 shaping term 不改变最优策略：

$$F(s, a, s') = \gamma \Phi(s') - \Phi(s)$$

其中 $\Phi$ 是 potential function。一个实用的 $\Phi(s)$ 例子是 distance-to-goal，提升 agent 在 early stages 的 guiding signal。

\begin{warningbox}{错误 shaping 的后果}
若 shaping 直接奖励“接近目标”而非“完成任务”，agent 可能提前停留在潜在 reward 高的状态，造成 reward hacking。这就是为什么 potential-based shaping 需严格依据 MDP structure。
\end{warningbox}

\subsection{本章小结}

Reward shaping 和 value stabilization 是保持 long-term training 稳定的核心。target smoothing、reward clipping 与 potential-based shaping 搭配使用可以有效避免 value explosion 与 reward hacking。

